\subsubsection{Gate Infidelity}\label{sec:gate-infidelity}

In the mathematical model of quantum computing, quantum gates are represented by ideal unitary matrices that transform quantum states exactly as predicted by theory.

\highspace
For example, the Hadamard gate is defined as:
\begin{equation*}
    H = \frac{1}{\sqrt{2}} \begin{bmatrix}
        1 & 1 \\
        1 & -1
    \end{bmatrix}
\end{equation*}
And when applied to the \ket{0} state, it always produces the \ket{+} state:
\begin{equation*}
    H \ket{0} = \frac{1}{\sqrt{2}} (\ket{0} + \ket{1}) = \ket{+}
\end{equation*}
In theory, this transformation is \textbf{exact}.

\highspace
Real quantum computers, however, do not manipulate abstract vectors and matrices. They manipulate physical systems through microwave pulses, laser pulses, electromagnetic fields, and control electronics. As a consequence, the operation implemented by the hardware is never perfectly identical to the ideal gate (as we have seen in the previous section).

\highspace
\begin{flushleft}
    \textcolor{DarkOrange3}{\faIcon{balance-scale} \textbf{Ideal Gate vs Physical Gate}}
\end{flushleft}
It is useful to distinguish between:
\begin{itemize}
    \item \textbf{Ideal gate}: the mathematical operation specified by the algorithm;
    \item \textbf{Physical gate}: the operation actually implemented by the hardware.
\end{itemize}
Ideally we would like:
\begin{equation*}
    U_{\text{physical}} = U_{\text{ideal}}
\end{equation*}
But in practice we have:
\begin{equation*}
    U_{\text{physical}} \approx U_{\text{ideal}}
\end{equation*}
But there is always a small discrepancy due to imperfections in the hardware. This \textbf{difference between the ideal and physical implementation} is called \definition{Gate Infidelity}.